\documentclass{article}
\usepackage[T1]{fontenc}
\usepackage{lmodern}
\usepackage{graphicx}
\usepackage{amsmath,amssymb}
\usepackage[a4paper,margin=2cm]{geometry}
\usepackage[absolute,overlay]{textpos}
\setlength{\TPHorizModule}{1mm}
\setlength{\TPVertModule}{1mm}
\usepackage[indonesian]{babel}
\usepackage{hyperref}
\setlength{\emergencystretch}{2em}
% unit: Halaman 37

\begin{document}

% === Halaman 37 (python/native) ===

37

def transposition\_cipher\_decrypt(ciphertext, k):

    jumlah\_baris = math.ceil(len(ciphertext) / k)

    jumlah\_sel\_kosong = (jumlah\_baris * k) - len(ciphertext)

    plaintext = [''] * jumlah\_baris     \# membuat sebuah list dengan panjang k, yang

                                        \# setiap elemennya berupa string kosong. 

    kolom = 0

    baris = 0

    for karakter in ciphertext:

        plaintext[baris] = plaintext[baris] + karakter

        baris = baris + 1

        if (baris == jumlah\_baris) or (baris == jumlah\_baris-1 and 

                                       kolom >= k - jumlah\_sel\_kosong):

            baris = 0

            kolom += 1

    return ''.join(plaintext)

\end{document}
